\section*{פתיחה: איפה אנחנו בקורס}
\begin{itemize}
  \item סיימנו דיפוזיה; כעת עוסקים במשוואת לפלס, גם בשיעור הזה וגם בשיעור הבא.
  \item לאחר מכן נעבור למשוואות אי-הומוגניות (הדברים ש``דחינו הצידה'').
  \item הדגש המרכזי: משפטי קיום ויחידות במשוואת לפלס --- אם \emph{מנחשים} פתרון שמקיים PDE ותנאי שפה, אז זהו \emph{הפתרון היחיד}.
\end{itemize}

\section{תזכורת: משוואת לפלס וחצי-מישור}
\subsection{הבעיה}
פותרים את משוואת לפלס
\begin{equation}
\Delta u = u_{xx}+u_{yy}=0
\end{equation}
בחצי-מישור
\[
y>0,\qquad x\in(-\infty,\infty),
\]
עם תנאי שפה על הציר:
\begin{equation}
u(x,0)=f(x).
\end{equation}

\subsection{למה לא במישור כולו?}
אם פותרים בכל המישור ומבקשים שבאינסוף הפוטנציאל ירד לקבוע (או ל-$0$), אז הפתרון הטריוויאלי
\[
u=\text{קבוע}
\]
מקיים את משוואת לפלס ולכן מתקבלת בעיה ``משעממת'' (אין מטענים במישור). בחצי-מישור ניתן לקבל תנאי שפה לא טריוויאליים שמייצרים פוטנציאל לא טריוויאלי בתחום.

\subsection{פתרון באמצעות טרנספורם פורייה (הרעיון)}
עושים טרנספורם פורייה ב-$x$:
\[
\mathcal{F}\{u\}(k,y)=R(k,y),
\]
כך ש-PDE הופכת ל-ODE:
\begin{equation}
\frac{d^2 R}{dy^2}-k^2 R=0.
\end{equation}
הפתרון הכללי הוא צירוף של אקספוננטים; בחצי-מישור בוחרים את החלק הדועך ב-$y$ כדי לא להתפוצץ באינסוף.

לאחר שמוצאים $R(k,y)$, חוזרים ל-$u$ בעזרת פורייה הפוך:
\begin{equation}
u(x,y)=\frac{1}{\sqrt{2\pi}}\int_{-\infty}^{\infty} R(k,y)e^{-ikx}\,dk.
\end{equation}

\subsection{פונקציית גרין וקונבולוציה}
בסוף מתקבל ניסוח של פתרון בתור קונבולוציה:
\begin{equation}
u(x,y)=\int_{-\infty}^{\infty} f(\xi)\,G(x-\xi,y)\,d\xi,
\end{equation}
כאשר $G$ היא פונקציית גרין של חצי-המישור (פתרון לפלס עם תנאי שפה דלתא).

\paragraph{משמעות פונקציית גרין.}
פונקציית גרין היא הפתרון כאשר תנאי השפה הוא
\[
u(x,0)=\delta(x-x_0),
\]
כלומר ``פוטנציאל דלתא'' על השפה. איזה התפלגות מטענים מייצרת תנאי כזה זו שאלה פיזיקלית נפרדת; למטרת פתרון PDE זה כלי פורמלי.

\subsection{השוואה לדיפוזיה: דעיכה}
\begin{itemize}
  \item בדיפוזיה פונקציית גרין דועכת גאוסיאנית (אקספוננציאלית).
  \item בלפלס בחצי-מישור מתקבלת דעיכה \emph{אלגברית} (למשל כמו $1/y$ במרחק גדול מהשפה, בהקשר הדו-ממדי שנדון).
\end{itemize}

\subsection{``מוד אחד'' על השפה}
אם תנאי השפה הוא מוד יחיד, למשל
\begin{equation}
u(x,0)=A\sin(kx)\quad(\text{או }A\cos(kx)),
\end{equation}
אז מתקבל מיידית פתרון מהצורה
\begin{equation}
u(x,y)=A\sin(kx)\,e^{-ky}.
\end{equation}
בדיקה מהירה: נגזרת שנייה ב-$x$ נותנת $-k^2u$ ונגזרת שנייה ב-$y$ נותנת $+k^2u$, ולכן $\Delta u=0$.

\subsection{יישום אינטואיטיבי: כלוב פראדיי}
תנאי שפה מחזורי (לא סינוס מושלם אלא מחזורי עם אורך גל טיפוסי $\lambda$) ניתן לפרק למודים, וכל מוד דועך אקספוננציאלית פנימה כמו $e^{-ky}$. לכן, ככל שנכנסים פנימה, הפוטנציאל והשדה דועכים ומתאפסים בקירוב --- מנגנון כלוב פראדיי.

\subsection{הערה: הפרדת משתנים מול פורייה בתחום אינסופי}
ניתן לפתור את חצי-המישור גם בהפרדת משתנים. ההבדל העקרוני:
\begin{itemize}
  \item בתחום סופי מתקבלים ערכים עצמיים דיסקרטיים $\Rightarrow$ סכומים.
  \item בתחום אינסופי מתקבל ספקטרום רציף $\Rightarrow$ אינטגרל על $k$.
\end{itemize}
מתברר שהשיטות נותנות את אותו מבנה פתרון.

\section{פרוסה (Strip): \(0<y<L\)}
\subsection{הבעיה}
\[
x\in(-\infty,\infty),\qquad 0<y<L,
\]
עם תנאי שפה
\begin{equation}
u(x,0)=f(x),\qquad u(x,L)=0,
\end{equation}
וכן דרישה לדעיכה ל-$0$ כש-$x\to\pm\infty$.

\subsection{פורייה ב-\(x\)}
שוב עושים טרנספורם פורייה ב-$x$ ומקבלים עבור $R(k,y)$:
\begin{equation}
R''-k^2R=0,\qquad
R(k,0)=F(k),\qquad
R(k,L)=0,
\end{equation}
כאשר $F(k)$ הוא טרנספורם פורייה של $f(x)$.

מפתרון שתי המשוואות מתקבל (בכתיב היפרבולי):
\begin{equation}
R(k,y)=F(k)\,\frac{\sinh\big(k(L-y)\big)}{\sinh(kL)}.
\end{equation}

\subsection{חזרה ל-\(u\) והגדרת פונקציית גרין}
הפתרון הוא פורייה הפוך:
\begin{equation}
u(x,y)=\frac{1}{\sqrt{2\pi}}\int_{-\infty}^{\infty}
F(k)\,\frac{\sinh\big(k(L-y)\big)}{\sinh(kL)}\,e^{-ikx}\,dk.
\end{equation}
מציבים $F(k)$ עצמו כאינטגרל על $\xi$:
\[
F(k)=\frac{1}{\sqrt{2\pi}}\int_{-\infty}^{\infty} f(\xi)e^{ik\xi}\,d\xi,
\]
ומקבלים קונבולוציה עם פונקציית גרין של פרוסה:
\begin{equation}
u(x,y)=\int_{-\infty}^{\infty} f(\xi)\,G_{\text{strip}}(x-\xi,y)\,d\xi,
\end{equation}
כאשר פורמלית
\begin{equation}
G_{\text{strip}}(x,y)=\frac{1}{2\pi}\int_{-\infty}^{\infty}
\frac{\sinh\big(k(L-y)\big)}{\sinh(kL)}\,e^{-ikx}\,dk.
\end{equation}
הערה מהשיעור: את האינטגרל הזה אפשר לפתור אנליטית אך מתקבל ביטוי מכוער (פונקציות מיוחדות); נומרית אפשר לחשב/לצייר בקלות.

\subsection{החלפת תנאי שפה}
אם במקום זאת
\[
u(x,0)=0,\qquad u(x,L)=g(x),
\]
אפשר לקבל את הפתרון על ידי החלפת המשתנה $y\mapsto L-y$ (כלומר להחליף בין $y$ ו-$L-y$ בביטויים).

\subsection{שני תנאי שפה יחד: סופרפוזיציה}
בבעיה הכללית בפרוסה:
\[
u(x,0)=f(x),\qquad u(x,L)=g(x),
\]
משתמשים בלינאריות:
\begin{itemize}
  \item $u_1$ פותר עם $f$ למטה ו-$0$ למעלה.
  \item $u_2$ פותר עם $0$ למטה ו-$g$ למעלה.
\end{itemize}
אזי
\begin{equation}
u=u_1+u_2
\end{equation}
מקיים גם את PDE וגם את שני תנאי השפה, ולפי משפט היחידות זהו הפתרון היחיד.

\section{תחום סופי מלבני: הפרדת משתנים}
\subsection{הבעיה (מלבן)}
תחום:
\[
0<x<L_x,\qquad 0<y<L_y,
\]
ותנאי שפה ``הכי פשוטים'':
\begin{equation}
u(x,0)=f(x),\qquad
u(x,L_y)=0,\qquad
u(0,y)=0,\qquad
u(L_x,y)=0.
\end{equation}

\subsection{הפרדת משתנים}
מחפשים
\[
u(x,y)=X(x)Y(y).
\]
הצבה בלפלס נותנת
\[
\frac{X''(x)}{X(x)}=-\frac{Y''(y)}{Y(y)}=-\lambda^2.
\]
הבחירה ב-$-\lambda^2$ נועדה לקבל ב-$x$ סינוסים/קוסינוסים שיכולים לקיים תנאי שפה הומוגניים בקצוות.

מקבלים:
\begin{align}
X''+\lambda^2X&=0,\\
Y''-\lambda^2Y&=0.
\end{align}

מתנאי שפה $X(0)=X(L_x)=0$ נקבל
\begin{equation}
\lambda_n=\frac{n\pi}{L_x},\qquad X_n(x)=\sin\left(\frac{n\pi x}{L_x}\right),\qquad n=1,2,\dots
\end{equation}

ל-$Y$ מתקבל פתרון אקספוננציאלי; עם תנאי $u(x,L_y)=0$ מתקבל מבנה היפרבולי:
\begin{equation}
Y_n(y)\propto \sinh\left(\lambda_n(L_y-y)\right),
\end{equation}
כך שב-$y=L_y$ אכן $Y_n=0$.

לכן פתרון כללי:
\begin{equation}
u(x,y)=\sum_{n=1}^{\infty} C_n\,
\sin\left(\frac{n\pi x}{L_x}\right)\,
\sinh\left(\frac{n\pi (L_y-y)}{L_x}\right).
\end{equation}

\subsection{קביעת המקדמים באמצעות תנאי השפה ב-\(y=0\)}
מציבים $y=0$:
\begin{equation}
f(x)=\sum_{n=1}^{\infty} C_n\,
\sin\left(\frac{n\pi x}{L_x}\right)\,
\sinh\left(\frac{n\pi L_y}{L_x}\right).
\end{equation}
מגדירים
\[
\alpha_n \coloneqq C_n\,\sinh\left(\frac{n\pi L_y}{L_x}\right),
\]
ואז $\alpha_n$ הם מקדמי טור סינוסים של $f$:
\begin{equation}
\alpha_n=\frac{2}{L_x}\int_0^{L_x} f(x)\,
\sin\left(\frac{n\pi x}{L_x}\right)\,dx.
\end{equation}
ולכן
\begin{equation}
C_n=
\frac{2}{L_x\,\sinh\left(\frac{n\pi L_y}{L_x}\right)}
\int_0^{L_x} f(x)\,\sin\left(\frac{n\pi x}{L_x}\right)\,dx.
\end{equation}

\subsection{החלפת תפקידים (תנאי שפה על צלע אנכית)}
אם תנאי השפה הלא-הומוגני הוא למשל על $x=0$ (והשאר $0$), אז מחליפים תפקידים: פונקציות עצמיות ב-$y$ ו-$x$ ``משועבד'' עם $\sinh$ ב-$x$.

\subsection{סופרפוזיציה במלבן}
אם יש תנאי שפה לא-הומוגניים על יותר מצלע אחת, מפרקים לסכום של בעיות פשוטות (כל פעם ``צלע אחת פעילה והשאר אפס'') ומחברים פתרונות.

\paragraph{אזהרה מרכזית.}
אם תנאי שפה אינם הומוגניים באופן שמונע הקמת בעיית שטורם--ליוביל נקייה, הפרדת משתנים ``ישירה'' נכשלת; לכן משתמשים בסופרפוזיציה ובטריקים נוספים כדי להפוך את הבעיה לאוסף בעיות עם תנאי שפה הומוגניים מתאימים.

\section{מעגל: מעבר לקואורדינטות פולריות}
\subsection{בעיה פנימית מול חיצונית}
\begin{itemize}
  \item \textbf{בעיה פנימית}: פותרים בתוך הדיסקה; אין מטענים בתוך התחום (אין סינגולריות), תנאי השפה על ההיקף.
  \item \textbf{בעיה חיצונית}: מטענים יכולים להיות בתוך הדיסקה, ומבקשים פתרון מחוץ לה.
\end{itemize}
בשיעור מתחילים בבעיה \textbf{פנימית}.

\subsection{למה לא לעבוד בקרטזיות?}
תנאי השפה על מעגל מערבב $x,y$:
\[
x^2+y^2=R^2,
\]
ולכן קשה לבצע הפרדת משתנים ב-$x,y$ תוך התאמה טבעית לתנאי השפה. עוברים לקואורדינטות פולריות/גליליות.

\subsection{לפלסיאן בפולריות (ללא תלות ב-\(z\))}
בהנחת אי-תלות ב-$z$:
\begin{equation}
\Delta u=
\frac{\partial^2 u}{\partial r^2}
+\frac{1}{r}\frac{\partial u}{\partial r}
+\frac{1}{r^2}\frac{\partial^2 u}{\partial \theta^2}
=0.
\end{equation}

\subsection{הפרדת משתנים \(u(r,\theta)=R(r)\Theta(\theta)\)}
מציבים ומחלקים ב-$R\Theta$:
\[
\frac{r^2R''+rR'}{R}=-\frac{\Theta''}{\Theta}=\lambda^2.
\]
הסימן נבחר כך ש-\(\Theta\) תהיה מחזורית.

\subsubsection*{החלק הזוויתי: תנאי מחזוריות}
המשוואה:
\[
\Theta''+\lambda^2\Theta=0
\]
עם תנאי מחזוריות:
\[
\Theta(\theta+2\pi)=\Theta(\theta).
\]
מכאן מתקבל \(\lambda=n\in\mathbb{Z}_{\ge 0}\) ופתרונות:
\begin{equation}
\Theta_n(\theta)=A_n\cos(n\theta)+B_n\sin(n\theta),\qquad n=0,1,2,\dots
\end{equation}
זהו בדיוק בסיס פורייה לפונקציות מחזוריות ב-\(\theta\) (המקרה הפרטי של שטורם--ליוביל עם תנאי שפה מחזוריים).

\paragraph{המקרה \(n=0\).}
אז \(\Theta''=0\), והמחזוריות מחייבת שהפתרון יהיה קבוע.

\subsubsection*{החלק הרדיאלי: משוואת כושי--אילר}
המשוואה הרדיאלית:
\begin{equation}
r^2R''+rR'-n^2R=0.
\end{equation}
מנחשים \(R(r)=r^\alpha\) ומקבלים \(\alpha=\pm n\), ולכן
\begin{equation}
R_n(r)=C_n r^n+D_n r^{-n},\qquad n\ge 1.
\end{equation}

\paragraph{בחירה פיזיקלית בבעיה פנימית.}
בדיסקה הפנימית $r=0$ שייך לתחום, ואין מטענים בתוך התחום, לכן $u$ לא יכול להתבדר ב-$r=0$.
מכאן של-$n\ge 1$ חייבים לזרוק את האיבר \(r^{-n}\), כלומר \(D_n=0\).

\paragraph{הערת מעבר.}
בסוף קטע השיעור (הטקסט שנמסר) מתחילים לשאול: מה קורה במקרה \(n=0\) בחלק הרדיאלי, ומהו פתרון המשוואה כאשר \(n=0\); המשך הטיפול מופיע בהמשך השיעור הבא.


